\chapter{Tecnologias}

Este capítulo apresenta as principais tecnologias empregadas no desenvolvimento
do aplicativo, organizadas de acordo com a camada da arquitetura em que
atuam: front-end, back-end e infraestrutura/persistência de dados.

\section{Front-end}

\subsection{React Native}

React Native é um \textit{framework} de código aberto, mantido pela Meta, que
permite o desenvolvimento de aplicativos móveis multiplataforma a partir de
uma única base de código escrita em JavaScript/TypeScript, utilizando
componentes que são traduzidos para elementos nativos tanto no iOS quanto no
Android \citep{reactnative:2026}.

A escolha do React Native se justifica pelo escopo reduzido da equipe de
desenvolvimento (um único desenvolvedor) e facilidade de conhecimento e
manutenção, sem trazer, para este projeto, benefícios que justifiquem esse
custo adicional. Além disso, o ecossistema maduro do React Native oferece
ampla disponibilidade de bibliotecas para funcionalidades já previstas no
aplicativo, como notificações locais e gráficos.

\subsection{Zustand}

Zustand é uma biblioteca leve de gerenciamento de estado para aplicações
React e React Native, que dispensa a estrutura verbosa de \textit{providers}
e \textit{reducers} tipicamente associada a bibliotecas como Redux
\citep{zustand:2026}.

Ela foi escolhida para gerenciar o estado efêmero da aplicação (por exemplo,
dados de formulários em preenchimento ou estados de navegação temporários)
por sua API minimalista e baixo \textit{overhead} de configuração, adequada a um
projeto de escopo individual em que a complexidade adicional de soluções
mais robustas de gerenciamento de estado não se justifica.

\section{Back-end}

\subsection{Express.js}

Express.js é um \textit{framework} minimalista para aplicações web em
Node.js, amplamente utilizado para a construção de APIs \textit{REST},
oferecendo uma camada fina de abstração sobre o roteamento de requisições HTTP e
\textit{middlewares} \citep{express:2026}.

Sua adoção se deve à simplicidade e à flexibilidade não opinativa do
\textit{framework}, que permite estruturar a API de acordo com as
necessidades específicas do projeto, sem impor convenções rígidas de
arquitetura, além de sua ampla documentação e maturidade no ecossistema
Node.js.

\subsection{Zod}

Zod é uma biblioteca de validação e definição de esquemas de dados para
TypeScript, que permite declarar formatos esperados de entrada (como corpos
de requisição) e inferir automaticamente os tipos estáticos correspondentes
\citep{zod:2026}.

Zod foi escolhido para validar os dados recebidos pela API antes de seu
processamento, garantindo que registros de humor, sono e demais entradas do
usuário estejam em conformidade com o formato esperado. A inferência
automática de tipos a partir dos esquemas evita a duplicação de definições de
tipo entre a camada de validação e o restante do código TypeScript do
back-end.

\subsection{Valkey e BullMQ}

Valkey é um banco de dados em memória, do tipo chave-valor, criado como
\textit{fork} de código aberto do Redis, mantido sob a governança da Linux
Foundation \citep{valkey:2026}. BullMQ é uma biblioteca para gerenciamento de
filas de tarefas em Node.js, construída sobre bancos de dados compatíveis com
o protocolo Redis \citep{bullmq:2026}.

A combinação das duas ferramentas foi adotada para o processamento paralelo
e assíncrono de tarefas que não precisam bloquear a resposta imediata da API,
como o disparo de notificações programadas e o processamento de indicadores
a partir do histórico de registros do usuário. A escolha do Valkey em
específico, e não do Redis diretamente, se deve à mudança de licenciamento do
Redis em 2024, que passou a restringir certos usos comerciais; o Valkey
mantém compatibilidade total de protocolo, permanecendo sob licença
permissiva de código aberto.

\section{Infraestrutura e Persistência de Dados}

\subsection{PostgreSQL}

PostgreSQL é um sistema gerenciador de banco de dados relacional, de código
aberto, reconhecido por sua conformidade com o padrão SQL, robustez e suporte
a tipos de dados avançados \citep{postgresql:2026}.

Sua escolha se justifica pela natureza estruturada e relacional dos dados do
domínio da aplicação (registros de humor associados a usuários, horários,
gatilhos e categorias), que se beneficia das garantias de integridade
referencial e das transações ACID oferecidas por um banco relacional
maduro, em detrimento de soluções não relacionais.

\subsection{Prisma}

Prisma é um \textit{Object-Relational Mapper} (ORM) para o ecossistema
Node.js/TypeScript, que gera automaticamente um cliente de acesso ao banco de
dados tipado a partir de um esquema declarativo, além de oferecer um sistema
de migrações versionadas \citep{prisma:2026}.

O Prisma foi escolhido por permitir a manutenção do esquema do banco de dados
de forma versionada e sincronizada com o código da aplicação, reduzindo a
chance de divergência entre a estrutura do banco e as entidades manipuladas
pelo back-end, além de prover checagem de tipos em tempo de compilação para
as consultas realizadas.
